\chapter{Conclusion: What Language Does Reality Force Us to Use?}

The inquiry began with the dispute between materialism and idealism, but the mathematical work forced a prior question into view. Before deciding whether matter, mind, form, or experience is fundamental, we must ask whether the descriptive language in which those categories are expressed preserves the distinctions reality itself requires. Ordinary language is indispensable for meaning and interpretation, yet words such as same, object, physical, mental, form, and change often compress relations that mathematics can separate. Mathematics in turn is not a single neutral language. A geometry, a dynamical system, a probability space, an invariant algebra, a quotient construction, and a Hilbert space preserve different distinctions. The problem is therefore not simply whether mathematics describes reality better than words, but which descriptive grammar loses the least structure while remaining no more complicated than the evidence requires.

Persistence exposed the problem first. A system can occupy genuinely different reduced states while instantiating the same form:

\[
\bar q_t\neq\bar q_s,
\]

while

\[
\Phi(\bar q_t)=\Phi(\bar q_s).
\]

The realization changes while the defining relation persists. Exact persistence and practical persistence then separate ideal invariance from bounded drift in real systems. The form map cannot be arbitrary; it must be justified by symmetry, causal organization, stable constraints, experimentally meaningful observables, or a declared identity criterion. Causal history may also be required where no instantaneous invariant is sufficient. The central insight is therefore not that relation is automatically metaphysically fundamental, but that identity in a dynamical universe need not be represented as equality of static state.

The oscillator models turn this insight into measurable dynamics. Collective rotation can be large while internal relational variation remains small. Stuart--Landau systems separate global phase, amplitude scaling, and internal deformation. These examples demonstrate that the whole can move without losing itself because movement of the whole and movement within the whole are mathematically different quantities. Coherence is therefore not the absence of change but preservation of defining relationship through change.

Platonic symmetry provides the clearest formal language of invariant form. Reynolds projectors measure the component of a field lying inside a finite-group fixed space, and character theory determines whether a particular symmetry is mathematically available at a chosen harmonic degree. The zeros in the multiplicity table provide exact exclusions; nonzero multiplicities provide only admissibility. Nonlinear dynamics must still determine stationarity, local stability, global preference, basin structure, and actualization frequency. This distinction between admissibility and realization becomes one of the deepest continuities between the mathematical and consciousness portions of the book.

The $\ell=5$ calculation then demonstrates that descriptive resolution itself has mathematical depth. The normalized quartic energy reduces exactly to a constant plus a positive multiple of $\|B_2(A)\|^2$, so its minimum set is the zero locus of $B_2$ on $S^{10}$. At the explicit regular point $A_\star$, the constraint differential has rank five. The local minimum manifold therefore has five tangent dimensions. Three are rotations of the same intrinsic form and two are genuine physical shape directions. Quartic order does not contain enough invariant information to distinguish those two directions energetically. Degree six contains four genuinely new invariant directions beyond lower-order products, and the explicit sextic $J_{4,8}$ has positive-definite intrinsic Hessian on the two-dimensional shape-flat plane. What quartic relation cannot see, sixth-order relation can distinguish.

The result must remain what it is. The flat directions are not a model of qualia and the sextic is not a theory of consciousness. Their importance is methodological. A restricted descriptive algebra can identify states that a richer algebra separates. A claim of completeness therefore requires evidence that the chosen descriptive language is sufficiently rich to separate the distinctions relevant to the question.

Wave dynamics adds a second dimension to the same argument. Geometry can classify form but cannot uniquely determine dynamics. One Laplacian can support diffusion, classical waves, Schrödinger evolution, higher-order unitary flow, and other laws. Thus $G\nRightarrow W$. Complex conjugation supplies an exact form/process separation: norm, density, symmetry coherence, and stabilizer can remain invariant while directed current reverses. The state changes in a manner invisible to the coarse form descriptors. Platonic duality and conjugation are nevertheless generally distinct operations. Mathematical precision therefore prevents suggestive words such as dual, opposite, and conjugate from being treated as interchangeable merely because they sound philosophically related.

These results motivate the first of the three governing principles: indistinguishability under a restricted descriptive algebra does not imply identity under every richer descriptive algebra. The second principle then prevents the argument from leaping from mathematical success to metaphysics: descriptive-language adequacy does not imply ontological sufficiency. The third applies the same discipline to consciousness: mathematical admissibility does not imply phenomenal occupancy.

These principles preserve rather than dissolve the materialism--idealism dispute. Sophisticated physicalism asserts that a complete physical relational description is sufficient. Idealism gives explanatory priority to experiential or intelligible structure. Relational dual-aspect realism proposes that exterior physical and interior experiential descriptions may be irreducible projections of one deeper actuality. None of these positions wins because its preferred language is more expressive. They become scientifically distinguishable only when the models make different predictions or impose different constraints on data.

The consciousness program therefore refuses to write $E$ directly into empirical equations. The latent $Z_\varphi$ represents converging phenomenological evidence and remains distinct from experience itself. Increasingly rich physical descriptions $P_{\mathrm{phys}}^{(k)}$ are tested for sufficiency by asking whether the conditional predictive contribution of $Z_\varphi$ disappears. Partial information decomposition separates redundant, unique, and synergistic contributions. Interventions test whether the proposed bridge has causal content. Cross-context transfer asks whether the same relation survives changes in task, subject, modality, and state. Complexity-matched physical latent models test whether an apparent relational advantage is merely additional model capacity. Explicit defeat conditions prevent the preferred ontology from becoming immune to evidence.

If systematic physical enrichment absorbs every phenomenologically relevant residual, relational physicalism gains support. If independently anchored phenomenological structure remains predictive and interventionally necessary after physical enrichment and complexity-controlled comparison, relational dual-aspect realism gains motivation. If experiential structure uniquely explains the physical observations in a way competing models cannot, idealism gains support. If the alternatives remain empirically equivalent, underdetermination is not an embarrassment to be hidden; it is the result the evidence warrants.

The final lesson is therefore not that one symbolic language should conquer all others. Words preserve semantic and experiential meaning that particular equations may not encode. Mathematics can preserve structural distinctions that ordinary language routinely collapses. A lower-order mathematical model may itself be too coarse and require a richer formalism. The correct task is to determine the mappings among reality, mathematical structure, conceptual meaning, and empirical observation without allowing any one mapping to pretend it is automatically identical with its source.

The language tells us how to ask the question. The evidence must decide what is real. The enduring principle that carries the argument from mathematics into philosophy remains simple: coherence is the persistence of defining relationship through variation. Its Platonic face is invariant form. Its wave face is transformation through form. Its epistemological face is descriptive resolution. Its consciousness boundary is admissibility versus occupancy. Its ontological test is whether the complete exterior physical description ultimately closes or whether reality continues to demand an irreducible interior relation that no improvement of exterior description can absorb.
